% ===========================================================================
%  问题四正文（由 main.tex 以 \input 引入）
%  表格由 03_代码/p4_tables.py 从 p4_results.json 自动生成：
%    p4_tables.tex          正文：两分支汇总 + 策略对照 + 组合比较
%    p4_tables_appendix.tex 附录：价格结构 + 预测精度 + 风险口径 + 校验（--split）
%  图由 03_代码/p4_figures.py 生成到 04_图/。
%  表号与图号由 LaTeX 自动编号，正文一律用 \ref，不写死数字。
% ===========================================================================

\section{问题四：实时波动电价下的两个分支}

\subsection{问题四新增的困难：价格从常数变成待预测量}

前三问的价格是一条已知的分时曲线：它在 $0{:}00$ 就已完全确定，优化器可以放心
地把它当作常数。问题四把价格换成附件 4 的\textbf{全年逐十分钟实际电价}，于是
价格本身成为\textbf{与负载、光伏并列的第三个不确定量}，而且它与后两者耦合：
电价高的时段往往是净负荷也高的时段。这带来三个必须分别处理的新问题：

\begin{itemize}
  \item \textbf{信息问题。}附件 4 已经把全年价格写在文件里，但"文件里有"不等于
        "$0{:}00$ 就知道"。决策时点只允许使用\textbf{当时已经实现的价格}；
        未来的价格必须预测，且必须能说清用的是什么信息。
  \item \textbf{耦合问题。}当电价与净负荷预测误差同向变动时，
        $\mathbb{E}[PR]\neq\mathbb{E}[P]\,\mathbb{E}[R]$，风险余量不能简单照搬
        问题二的等权分位数。
  \item \textbf{口径问题。}两个交付文件
        （\texttt{result4-2.xlsx} 与 \texttt{result4-3.xlsx}）对应两类策略权限，
        它们必须\textbf{在同一条实际价格路径上各自独立回放}才能比较，否则比出来
        的差额里会同时混入"策略变了"和"计价标准变了"两个因素。
\end{itemize}

本节余下的安排是：先分解附件 4 的价格结构（第 \ref{sec:p4-struct} 节），据此
导出可用的价格预测（第 \ref{sec:p4-forecast} 节），再处理价格的联合风险
（第 \ref{sec:p4-risk} 节），然后把两个分支的模型写清楚
（第 \ref{sec:p4-model} 节），最后给出结果、策略对照与组合比较。

\subsection{附件 4 电价的结构：与附件 1 同水平、多一个星期维度}
\label{sec:p4-struct}

附件 4 共 $365\times144=52{,}560$ 个价格，$0.0076$--$1.7936$~元/kWh，均值
$0.7662$~元/kWh，标准差 $0.3407$~元/kWh，日中位峰谷比 $4.33$。它有两个来源
完全不同的变异：\textbf{日内形状}（傍晚贵、凌晨便宜）与\textbf{日间水平}。
若用 $365$ 天的分段均值刻画日内形状（这用到了全样本信息，是\textbf{事后口径}，
只能当作可预测性的上界），\textbf{单一日内形状即可解释 $84.6\%$ 的方差}，
而日间水平的离散度——逐日均值的标准差除以均价——只有 $0.135$，这是一个
\textbf{无量纲的变异系数}（对应标准差 $0.1036$~元/kWh）。两者合起来说明
"哪一天"的影响远小于"一天里的哪个时段"。真正的问题是：日间那部分变异里，
有没有\textbf{可利用的规律}？

把附件 4 记为 $P$、附件 1 的固定曲线记为 $F$，做分解 $P=F+R$，得到四个断言
（数值见表~\ref{tab:p4-struct}，图~\ref{fig:p4-struct} 左）：

\begin{enumerate}
  \item \textbf{水平相同。}$P$ 的 $365$ 天逐时段均值与 $F$ 的最大偏差只有
        $5.15\times10^{-5}$~元/kWh，两者年均值分别为 $0.766198$（$P$，附件 4）
        与 $0.766197$（$F$，附件 1），相差不足 $10^{-6}$~元/kWh。
        这条最关键：它决定了后文所有对照实验的读法——既然价格水平没变，
        任何费用差异都\textbf{只来自波动本身}，不掺杂涨价或降价。
  \item \textbf{残差与骨架正交。}$\mathrm{corr}(R,F)=-1.5\times10^{-5}$。
        故 $R$ 就是"波动"本身，不含任何原先已知的分时结构。
  \item \textbf{$R$ 有星期结构。}把 $R$ 按星期几分组求均值，周一至周四与周日
        都在 $+0.053$~元/kWh 附近，\textbf{周五与周六却是 $-0.137$}~元/kWh。
        两个水平之间差 $0.1913$~元/kWh，相当于均价的 $25.0\%$。这是残差里
        \textbf{唯一}可被$0{:}00$ 时刻利用的规律。
  \item \textbf{噪声占多数。}扣除星期效应后，$R$ 的方差只下降 $41.7\%$，
        其余 $58.3\%$ 是无规律的日内噪声，无论用什么模型都预测不掉。
\end{enumerate}

由此可以给出价格预测的\textbf{推导}而不是\textbf{设想的}形式。记 $\bar R_{w}$
为星期 $w$ 的残差均值（由历史估计），则对 $n$ 天的第 $t$ 段，
$$
\boxed{\;
\widehat p_{n,t}=\frac{\sum_{i\in\mathcal L}w_i\,p_{n-i,t}}
{\sum_{i\in\mathcal L}w_i}
\;},\qquad \mathcal L=\{7,14,21\},\quad w_1\ge w_2\ge w_3>0,
\label{eq:p4-price}
$$
即\textbf{取最近三个同星期几的价格加权平均}。它的依据是恒等式
"同星期加权 $=$ 附件 1 曲线 $+$ 同星期残差加权"：附件 1 曲线提供日内形状，
而最近三个同星期几的实测价格则估计那个\textbf{缓慢漂移的星期偏移}。该恒等式
在 $334$ 天上的最大数值残差为 $4.44\times10^{-16}$~元/kWh，即它是同义改写
而非近似。权重取递减（越近权重越大）是因为星期一层的偏移会漂移，越近的样本
越可信。\textbf{具体取哪一组权重不在评价区间上决定}，见
\ref{sec:p4-forecast} 节的选型程序。

\begin{figure}[htbp]
  \centering
  \includegraphics[width=.88\textwidth]{p4_fig1_price_structure.png}
  \caption{附件 4 电价的结构分解：左为实际电价与附件 1 曲线（横轴为一周），
           右为残差 $R$ 的星期均值}
  \label{fig:p4-struct}
\end{figure}

\subsection{价格预测：预热期选型、评价区间内冻结}
\label{sec:p4-forecast}

框架 3.1 节要求"先做简单、可追踪的预测，并明确误差的评价方式"，只有在给出
明确理由时才引入更复杂的模型；框架第八节第 2 条又要求"选参同样只用过去
数据"。这两条合起来意味着：\textbf{预测器的超参数不能在评价区间上挑}。
否则"用 $2$--$12$ 月的数据选定权重与 $\beta$、再在 $2$--$12$ 月上报出精度"
就等于参考了答案再选方法，报出的误差必然是乐观的。

本文因此把样本切成两段，选型与评价\textbf{严格分开}：\textbf{选型期}为
$2025$ 年 $1$ 月（第 $1$--$30$ 天，$30$ 天、$4{,}320$ 个时段），权重 $w$ 与
日内修正系数 $\beta$ \textbf{只在这里选一次}；\textbf{评价期}为 $2$ 月 $1$ 日
--$12$ 月 $31$ 日（$334$ 天、$48{,}096$ 个时段），选型结果\textbf{原样冻结}，
选型期一过就不再改动任何一个参数。

选型程序是：把常见候选（昨天同时段、上周同一天、近 $7/14/28$ 天滚动均值、
同星期等权、两组同星期加权）在选型期上比平均绝对误差，取最优者。结果见
附录表~\ref{tab:p4-select}：\textbf{同星期加权 $0.5/0.3/0.2$} 以
$0.0486$~元/kWh 胜出，其后是同星期等权 $1/3$（$0.0487$）与加权
$0.6/0.3/0.1$（$0.0488$），三者几乎并列；而所有\textbf{跨星期滚动均值}候选
与"昨天同时段"（$0.0847$--$0.0861$）都差出近一倍——把不同星期几混在一起平均，
恰好把电价里唯一可利用的 $0.1913$~元/kWh 星期偏移\textbf{抹平}了。该表同时把
选型期切成三个子窗口重跑，\textbf{第一名都是 $0.5/0.3/0.2$}，排除"$30$ 天窗口
的退化期造成的偶然"。

选型一旦跳出评价区间，\textbf{代价也就看得见了}。把同样的候选放到评价区间
上重排（表~\ref{tab:p4-forecast}），第一名变成 $0.6/0.3/0.1$（$0.0441$），
而我们冻结使用的 $0.5/0.3/0.2$ 是 $0.0446$——相差 $0.0005$~元/kWh，约
$1.1\%$。本文没有取这个便宜：这 $0.0005$ 正是"参考答案选方法"换来的虚高
精度。表~\ref{tab:p4-forecast} 因此不是选型依据，而是\textbf{样本外复核}。

\textbf{日内修正。}$0{:}00$ 定下全天预测后，随着当天价格逐段实现，可以用已实现
时段对剩余时段做一次日内修正。取修正量
$$
\widehat p_{n,t}\leftarrow\widehat p_{n,t}\,
\exp\!\Big(\beta\cdot\frac{1}{|\mathcal T_{<}|}\sum_{s\in\mathcal T_{<}}
\ln\frac{p_{n,s}^{\mathrm{act}}}{\widehat p_{n,s}}\Big),\quad t\ge t_0,
$$
即用已完成时段的\textbf{平均对数偏差}外推剩余时段。$\beta$ 同样\textbf{只在
选型期扫描一次}，且\textbf{逐发布时刻单独选}（各时刻可用于估计偏差的已完成
样本数不同：$6{:}00$ 只有 $36$ 段、$18{:}00$ 有 $108$ 段）。$1$ 月扫描得
$\beta=(1.0,\,0.75,\,0.5)$，本文全程冻结于此。若改到评价区间上扫
（表~\ref{tab:p4-beta}），$6{:}00$ 版会指向 $\beta=0.5$、比冻结值好看
$0.0009$~元/kWh——那同样只是区间自己的偏好，本文不取其便。

\textbf{一个必须说明的可比性问题。}第 $h$ 个版本只能修正它发布之后的时段，
因此四个版本的"平均绝对误差"覆盖的时段\textbf{互不相同}（$0{:}00$ 版覆盖全天
$144$ 段，$18{:}00$ 版只覆盖最后 $36$ 段），\textbf{直接比大小是错的}。
表~\ref{tab:p4-skill} 因此另设"共同时段"一列，把四版都限制在
$18{:}00$--$24{:}00$ 上比较：$6{:}00$ 版最好（$0.0389$），其后是 $18{:}00$
版（$0.0403$）与 $12{:}00$ 版（$0.0409$），$0{:}00$ 版最差（$0.0428$）。
三版更新都优于不修正的 $0{:}00$ 版，但\textbf{并非"发布越晚越准"}：$12{:}00$
版反比 $18{:}00$ 版差 $0.0006$~元/kWh。这一倒挂的成因已查清，见附录
\ref{app:p4-inversion} 节，此处不粉饰。作为参照，完全不用日内形状信息
（全天套用均价）的误差为 $0.2774$~元/kWh，是四版的六倍以上。

\subsection{联合风险：高价恰逢缺电时不能把两个期望相乘}
\label{sec:p4-risk}

问题二的风险余量对净负荷预测误差 $\varepsilon$ 取等权分位数。在波动电价下，
真正决定损失的是乘积 $p\cdot\varepsilon$，而
$$
\mathbb{E}[PR]=\mathbb{E}[P]\,\mathbb{E}[R]+\mathrm{Cov}(P,R).
$$
\textbf{若 $\mathrm{Cov}(P,R)\neq0$，则不能把两个期望直接相乘。}因此第一步不是
假设，而是\textbf{检查协方差}，结果见表~\ref{tab:p4-cov}：合并口径的相关系数
只有 $+0.0826$，但\textbf{把时段固定住}以后，$144$ 个时段的逐时段相关系数均值
为 $+0.206$（范围 $+0.083$--$+0.401$，\textbf{$100.0\%$ 的时段为正}）。合并口径
之所以低，是被"时段之间"的电价差异主导了；同一时段内，$\varepsilon$ 与 $p$
确实是同向的。按电价高低分组看得更直接：电价最高的 $20\%$ 时段里
$\varepsilon$ 平均为 $+38.5$~kW，最低的 $20\%$ 里是 $-72.7$~kW，两端相差
$111.2$~kW。

结论有两条，都很具体：其一，$\mathrm{Cov}(P,R)>0$ 是\textbf{同一时段内}的性质，
故风险余量必须\textbf{逐时段}计算，不能对全年做一次汇总——这正是问题二
逐时段求分位数、而不是对全天求和后再加余量的原因；其二，既然
$\mathrm{Cov}>0$，忽略它就会\textbf{低估}紧急购电的风险，不能用简单的
$\mathbb{E}[P]\mathbb{E}[R]$ 近似。

\textbf{价格加权分位数。}据此把风险余量从等权分位数改为\textbf{以实际电价为
权重的经验分位数}：设第 $n$ 天第 $t$ 段的历史误差样本为
$\{\varepsilon_{i,t}\}_{i<n}$，对应权重取\textbf{同一目标时段}的实际电价
$\{p_{i,t}\}_{i<n}$，则
$$
\widetilde N_{n,t}=\widehat N_{n,t}+Q^{P}_{\alpha}\big(\varepsilon_{\cdot,t}\big),
\qquad
Q^{P}_{\alpha}=\inf\Big\{x:\
\frac{\sum_{i}p_{i,t}\mathbf 1[\varepsilon_{i,t}\le x]}{\sum_i p_{i,t}}\ge
\alpha\Big\}.
$$
它就是\textbf{加权经验分布的下确界}：把样本按误差升序排列，取\textbf{累积权重
比例首次达到 $\alpha$ 的那个样本值}。取样本值而不在两个样本之间插值，是因为
插值点上的加权覆盖率并不等于 $\alpha$（样本 $[0,10]$ 等权、$\alpha=0.8$ 时插值
给 $6$，而该点的加权覆盖率只有 $0.5$）；下确界定义才是临界比规则所要求的分位数，
且权重全相等时它精确退化为经验分布分位数。当权重和为零（样本缺失）时退回普通
分位数，避免除零。这样"贵的时段"在分位数里占更大权重，余量被推向\textbf{代价
更高}的时段。

表~\ref{tab:p4-risk} 给出两种口径的实测差。以\textbf{同约定}的等权下确界为
基准，$\alpha=0.6$ 时价格加权高 $3.99$~kWh/时段（$+37.2\%$），$\alpha=0.7$ 时
高 $3.02$~kWh（$+11.4\%$），$\alpha=0.8$ 时高 $1.51$~kWh（$+3.3\%$）；该表同时
列出以问题二线性插值口径为基准的对照（$+3.00/+3.57/+4.31$~kWh），两者之差纯粹
来自取值约定。无论取哪个基准，抬高的绝对量都在 $3\sim4$~kWh/时段量级，相对于约
$110$~kWh/时段的风险净负荷只有百分之几，\textbf{并不大}——原因是同一时段的价格
在日与日之间变动有限，权重差异没有合并口径暗示的那么悬殊。本文如实报告这一点：
这是一项\textbf{方向正确、幅度温和}的改进，而不是决定性的修正。之所以保留它，
是因为它的实现成本为零，且方向与上面实测出的协方差符号一致。需要说明的是，
上表衡量的是\textbf{风险余量本身}的差别，不是最终账单的差别；余量抬高
$3\sim4$~kWh/时段后，经过日计划 MILP 与逐段执行两层传导，落到总费用上的差异
会进一步缩小，因此本文不声称该改动带来了可辨认的费用下降，只声称它在方向上
更符合 $\mathrm{Cov}(P,R)>0$ 的事实。

\subsection{计费口径与两个分支的改造}
\label{sec:p4-model}

\subsubsection{结算原则：决策用当时的信息，一切按实际价格结算}

所有分支遵守同一条结算原则，它也是问题四全部代码的统一约束：

\begin{center}
\footnotesize
\begin{tabular}{lll}
  \toprule
  环节 & 使用什么价格 & 说明 \\
  \midrule
  \textbf{决策}（定 $g^0$、定 $x$、定储能轨迹）
    & 当时\textbf{可获得}的价格信息 & 决定"买多少、何时买" \\
  \textbf{结算}（算账）
    & 附件 4 的\textbf{实际价格} & 决定"花了多少钱" \\
  \bottomrule
\end{tabular}
\end{center}

因此，$4$-2 与 $4$-3 的全部费用系数（计划费、调整费、紧急费）都乘\textbf{对应
交付时段的实际电价}：
$$
C=\underbrace{\sum_t p_t^{\mathrm{act}}\,g^0_t}_{\text{计划费}}
+\underbrace{1.5\sum_t p_t^{\mathrm{act}}\big(x_t-g^0_t\big)_+}_{\text{调整费（上调）}}
+\underbrace{5\sum_t p_t^{\mathrm{act}}r_t}_{\text{紧急费}} ,
\qquad (y)_+=\max(y,0).
$$
这里\textbf{计划费的基数是初始计划 $g^0_t$，不是最终生效量 $x_t$}，这一点必须
写清：$1.5p_t$ 是上调部分的\textbf{总价}，而不是"先按 $p_t$ 付一遍、再按
$1.5p_t$ 付一遍"成为 $2.5p_t$。若把首项写成 $\sum_t p_t^{\mathrm{act}}x_t$，
由于 $x_t=g^0_t+a_t$，上调部分就被计了 $1+1.5=2.5$ 倍，与题目"超出部分的电价是
交易时刻电价的 $1.5$ 倍"直接矛盾。本文代码实现的是上式（按 $\sum_t p^{\mathrm{act}}g^0_t$
结算计划费，上调增量单独按 $1.5p^{\mathrm{act}}$ 计），问题三第 7.3 节已对此
作过同样的说明。
主计费口径沿用问题三的结论：原计划已按 $p^{\mathrm{act}}$ 全额付费，故保留原量
并把多余电量无偿弃用，比减少购电再付违约费\textbf{至少不贵}，因而令
$x_t\ge g^0_t$、调整费退化为上式的单项。第 \ref{sec:p4-result} 节会验证
"下调量在浮点零的意义下恒为零"，与这一理论结论一致。退款口径作为独立敏感性
单独保留，不与主口径混合。

\textbf{目标值与账单必须分开。}$0{:}00$ 的优化器最小化的是\textbf{按预测价格}
计算的费用 $C^{\mathrm{plan}}=\sum_t\widehat p_t x_t+\cdots$，而账本记录的是
$C^{\mathrm{act}}=\sum_t p_t^{\mathrm{act}}x_t+\cdots$。两者不相等，差额
$C^{\mathrm{act}}-C^{\mathrm{plan}}$ 就是"预测觉得自己省了多少"与"账本上真的
省了多少"的差。所有内部明细表同时输出\textbf{预测电价、实际电价、预估计划费、
实际计划费}四列，使这一差额可以在逐时段粒度上定位，而不只在总数上体现。

\subsubsection{分支 4-2：计划全天冻结，价格只在 $0{:}00$ 用一次}

4-2 对应交付文件 \texttt{result4-2.xlsx}：全天 $144$ 段的购电量在 $0{:}00$
一次定下，\textbf{日内不得修改}。相对问题二的日计划 MILP，只替换两处输入：

\begin{enumerate}
  \item 目标函数中的价格系数由附件 1 的固定曲线换成 $0{:}00$ 的\textbf{价格
        预测} $\widehat p^{0}$，即 $\min\sum_t\widehat p^{0}_t g_t+\lambda\xi$；
  \item 风险净负荷由等权分位数换成第 \ref{sec:p4-risk} 节的价格加权分位数
        $\widetilde N_t=\widehat N_t+Q^{P}_{\alpha}(\varepsilon_{\cdot,t})$。
\end{enumerate}

负载与光伏预测\textbf{仍然只用历史外推}（问题二口径），$4$-2 不使用附件 3 的
光伏预报——这是它与 $4$-3 的关键区别，也保证了两个分支的差异只来自"是否允许
调整、是否使用新预报"。日内执行完全复用问题二的因果规则：按当前实际负载与
光伏计算盈余、按参考轨迹构造储备线、充放电与紧急补购逐段决定。\textbf{执行层
唯一的外部常规购电来源是 $0{:}00$ 定下的 $g$}，日内不存在第二次常规购电，
这正是"计划冻结"在代码中的实际含义，也是第 \ref{sec:p4-result} 节校验它的
方式。

\subsubsection{分支 4-3：价格与光伏预报共同驱动的滚动调整}

4-3 对应 \texttt{result4-3.xlsx}，它把问题三的发布与提交规则\textbf{原样保留}，
只把价格换成波动电价：沿用交付块策略（$0{:}00$ 提交全天，$6{:}00$、$12{:}00$、
$18{:}00$ 各只提交自己那一个六小时块）、沿用三个购电量的定义
（$g^0$、$x$、$r$）、沿用问题二的实际执行规则。唯一的、也是本质的改变是：
\textbf{调整模型的目标函数用当时的预测价格}。设 $\mathcal T_\tau$ 为发布时刻
$\tau$ 到当日 $24{:}00$ 的尚未执行时段，
$$
\min\ \sum_{t\in\mathcal T_\tau}1.5\,\widehat p^{(\tau)}_t
\big(y_t-g^0_t\big)
+\sum_{t\in\mathcal T_\tau}5\,\widehat p^{(\tau)}_t \bar r_t+\lambda\xi ,
$$
其中 $\widehat p^{(\tau)}$ 是\textbf{$\tau$ 时刻可用版本}的价格预测：
$0{:}00$ 用 $0{:}00$ 版，$6{:}00$ 用第 \ref{sec:p4-forecast} 节修正后的
$6{:}00$ 版（对剩余时段）。因此 4-3 相对 4-2 多出的能力是\textbf{双重的}：
它在 $6{:}00$ 之后既能用新发布的光伏预报，也能用当天\textbf{已经实现的
价格}修正剩余时段的电价预测——而 4-2 两者都没有。

\textbf{信息边界由数据结构保证，而非靠调用约定。}第 $\tau$ 次求解只读取
$F^{(\tau)}$ 与 $\widehat p^{(\tau)}$ 两张表，它们都只含 $\tau$ 时刻已经发布／
已经实现的信息。另外，初始计划按 $(\text{价格模式},\text{计费口径},
\text{预报组合})$ 三元组分别缓存，\textbf{不同价格模式绝不能互相复用计划}，
否则"用固定价定的计划"会漏进"波动价策略"的回放里。这一条在代码中用缓存键
强制保证，并在第 \ref{sec:p4-result} 节校验。

\textbf{执行与跨日衔接}与问题三完全一致：$b_t=x_t+v_t^{\mathrm{act}}-
\ell_t^{\mathrm{act}}$，$c_t=\min\{(b_t)_+,M,(10800-E_{t-1})/\eta_c\}$，
$d_t=\min\{(-b_t)_+,M,\eta_d(E_{t-1}-R_t)_+\}$，$r_t=(-b_t-d_t)_+$，逐段守恒
$x_t+r_t+v_t^{\mathrm{act}}+d_t=\ell_t^{\mathrm{act}}+c_t+w_t$；当日日末电量直接
作为次日初值，两个分支\textbf{各自独立保存}实际状态，$1$ 月为预热期。

\subsection{求解结果与校验}
\label{sec:p4-result}

\subsubsection{两个交付文件的结果}

以 $2025$ 年 $1$ 月为预热期、$2$ 月 $1$ 日--$12$ 月 $31$ 日共 $334$ 天为正式
区间，两个分支各自独立标定参数，并在\textbf{同一条附件 4 实际价格路径}上回放，
结果见表~\ref{tab:p4-branch}。\textbf{result4-2} 在计划冻结的约束下全年合计
购电费 $14740813.21$~元；\textbf{result4-3} 允许在 $6{:}00$、$12{:}00$、
$18{:}00$ 三个时刻调整计划，全年合计 $14402403.87$~元，比 4-2 少
$338409.34$~元（$-2.30\%$）。

两个分支的费用结构差异比总费用更能说明问题。4-2 无法调整，缺口只能按 $5$ 倍
电价紧急补购，紧急费在合计中占比较高；4-3 用\textbf{两档更便宜的手段}替代了
它——一是在 $6{:}00$ 之后按 $1.5$ 倍电价追加（调整费），二是凭已实现的价格与
新发布的光伏预报\textbf{提前}把电量买在便宜时段。由于 $1.5p<5p$ 且有相当一部分
缺口被提前识别，这一置换在总体上有利。需要提醒的是，\textbf{两个分支的费用
不可相加}：它们是同一道题同一段价格路径下的两种规则。

期末储电量也必须一并报告：若某一分支靠在年末放空储能"换"出更低的账面费用，
那是把成本转嫁给了下一期。表~\ref{tab:p4-branch} 给出两者的期初、期末储电量，
二者处在同一量级，故上面的费用比较是公平的。

\textbf{标定结果本身也说明了价格波动改变了什么。}$\lambda$ 由 $1$ 月的 $180$ 组
$(\alpha,\rho,\lambda)$ 联合标定选出后\textbf{整年冻结}：4-2 选中 $\lambda=0.80$、
4-3 选中 $\lambda=0.50$，都远高于同一内核在问题二里选出的 $\lambda=0.00$
——日末储备这只"影子价格"，只有在价格逐日波动时才真正值钱。

$4$-$2$ 的 $(\alpha,\rho)$ 在 $8$ 个标定窗口里\textbf{真实摆动}，每个窗口的
胜出者都比次优者低 $26.9$--$3112.5$ 元，是一次有区分度的选择；方向也符合预期
——计划在 $0{:}00$ 冻结，"买了不能退、少了要按 $5$ 倍补"是不对称风险，最优
回应是\textbf{系统性地多买一点}，故 $\alpha$ 偏向高分位。

$4$-$3$ 的这一项\textbf{不能作同样解读}：$8$ 个窗口\textbf{全部}选中搜索网格的
左下角 $\alpha=0.70,\rho=0.5$，且网格内窗口费用对两者\textbf{单调递增}（$392$
次相邻比较中 $390$ 次成立），说明真实最优\textbf{不在网格之内}；其中 $4$ 个
窗口的胜出者与次优者只差 $0.19$--$25.12$ 元，相对窗口费用（约
$2.4\times10^{6}$ 元）是 $10^{-5}$ 量级，属于\textbf{数值平局}。因此这条记录
\textbf{既不能读成"$4$-$3$ 偏好低分位"，也不能读成"$4$-$3$ 的安全垫比
$4$-$2$ 薄"}，它反映的是\textbf{边界解加取点任意性}，本文把它列入模型边界
（第 \ref{sec:p4-limit} 节）而非结论。逐窗口的完整标定记录见附录。

两分支的标定差异因此只能谨慎陈述：4-2 的高 $\alpha$ 可以解释，4-3 的参数选择
则\textbf{不具可解释性}；4-3 的全年计划量 $21344253$~kWh 低于 4-2 的
$21705435$~kWh，但其中有多少来自"决策可修正"、多少来自"参数恰好停在角上"，
本文\textbf{没有拆开}。

\subsubsection{逐日校验}

对两个分支的 $334$ 个交付日逐日执行六项检查，结果见表~\ref{tab:p4-check}，
\textbf{未通过天数全部为零}。六项检查覆盖信息边界（求解读取的预报其发布时刻均
不晚于求解时刻）、计划冻结（4-2 的 $g_t$ 恒等于 $0{:}00$ 定下的 $g_t^0$）、
调整权限（4-3 恒有 $x_t\ge g_t^0$）、物理可行性（供需平衡残差在机器精度量级、
无同时充放电）、状态累计（储电量序列与记录值一致，日末电量严格作为次日初值）
与费用一致性（按预测价算的"预估费用"与按\textbf{实际}价格结算的费用分别复算，
均与记录值相符）——这正是第 \ref{sec:p4-model} 节的原则：\textbf{决策用当时的
信息，结算一律用实际价格}。逐条明细见附录第 \ref{app:checks} 节。

\subsubsection{三种价格信息条件下的策略对照}
\label{sec:p4-strategy}

要回答"适应波动电价值多少钱"，必须把\textbf{决策时能看到的价格信息}作为唯一
的变量：三种策略在同一套负载、光伏、储能与执行规则下运行，也都在\textbf{同一
条实际价格路径上结算}，只有制定计划时所依据的价格不同，因此费用可以直接相减。
三者分别是：

\begin{itemize}
  \item \textbf{波动电价预测策略}（主策略）：$0{:}00$ 用第
        \ref{sec:p4-forecast} 节的价格预测制定全天计划；4-3 在三个发布时刻用
        当天已实现的价格修正剩余时段。
  \item \textbf{固定电价参考策略}：直接拿附件 1 的分时电价曲线当全天价格决策，
        但\textbf{同样按附件 4 的实际价格结算}，模拟"没注意到电价已放开"的情形。
  \item \textbf{未来价格已知参考}：假设当天价格在 $0{:}00$ 已全部公布，是
        \textbf{信息增强上界}；题目未提供这一信息条件，故\textbf{不作为主策略、
        也不计入 $\Delta J$}，只用来度量"价格预测还有多少改进空间"。
\end{itemize}

记 $\Delta J=J(\text{固定电价})-J(\text{波动电价预测})$，即适应波动电价
值多少钱。三种策略\textbf{各自独立标定}参数，用的是各自信息集下的历史；否则
"用波动电价调出来的参数"会让固定电价策略凭空受益，比较就掺进了调参差异。
结果见表~\ref{tab:p4-strategy} 与图~\ref{fig:p4-strategy}。两个分支上
$\Delta J$ \textbf{都为正}：4-2 为 $20139$~元（$+0.14\%$），4-3 为
$32170$~元（$+0.22\%$）。方向符合直觉——按预测的波动价格安排购电，比照搬
目录电价曲线更省——但\textbf{幅度很小}，两者都不到全年费用的千分之四。
这一点必须如实说明，不能因为结论"对我们有利"就只报方向。

为什么适应波动价格值不了更多的钱？把两个参考策略拆开看就清楚了。

\textbf{其一，附件 1 曲线已经包含了绝大部分可预测的形状信息。}第
\ref{sec:p4-struct} 节已量化过：单一日内形状解释了价格方差的 $84.6\%$，而固定
电价策略用的正是这条形状曲线——它并不是什么都不懂的对手，已经知道傍晚贵、
凌晨便宜。真正需要预测的只剩残差 $R$，其中可被利用的只有星期偏移（$41.7\%$）。

\textbf{其二，两个策略的费用结构差异恰好印证了这一点。}在 4-3 上，$32170$
元的净优势\textbf{几乎全部来自计划费}：波动电价预测策略的计划费比固定电价参考
低 $33786$~元（$13778495$ 对 $13812281$），因为它看得出当天哪些时段便宜、把
电量排到低位；这份"敢买"也让计划离实际负荷更远，事后多付 $1778$ 元调整费、
少付 $162$ 元紧急费，合计只多 $1616$~元，可以忽略。换言之，$\Delta J$ 的符号
与量级\textbf{由计划费一项决定}。固定电价参考看不出价格的高低差，只能保守地
少买、再被动地补。

\textbf{其三，剩下的差距就是价格预测的改进空间。}"未来价格已知参考"在两个
分支上分别为 $14646668$ 与 $14293927$~元，比波动电价预测策略再低
$94145$~元（$0.64\%$）与 $108477$~元（$0.75\%$）。这就是\textbf{当天的
价格若能完全预知}所能节省的上限——它给出了价格预测还有多少可挖掘的余量，
也界定了本文 $0.0446$~元/kWh 的预测误差对应多大的经济代价（该精度即第
\ref{sec:p4-forecast} 节冻结预测器在评价区间上的样本外 MAE）。这一上界\textbf{不
计入 $\Delta J$}：题目并未提供"当天价格提前全部公布"这一信息条件。三种策略
的费用分解图见附录图~\ref{fig:p4-strategy}。

综合来看，本题的结论是：\textbf{在本文的价格序列下，"适应波动电价"确实有
价值，但价值有限}。原因不是策略不好，而是\textbf{波动里可预测的成分本来就少}
——价格与目录曲线的差异绝大部分是随机的星期与日内噪声。这也提示，若现实中的
价格波动包含更强的可预测结构，$\Delta J$ 会显著大于此处的量级（$2.0$--$3.2$ 万元）。

\subsubsection{八种预报组合}

把 4-3 的三个额外预报时刻按 $2^3=8$ 种方式取出，每种组合\textbf{独立重新
标定}后完整回放全年，可以回答在\textbf{价格本身也要预测}的条件下"还需要哪些
时刻的预报"，结果见附录表~\ref{tab:p4-combo}。全部弃用预报时费用为
$15022901$~元，比主策略（三版全用，$14402404$~元）高 $620497$~元（$+4.13\%$），
即三版预报合起来值 $62.05$~万元。但\textbf{三个时刻并不等价}：单独多订一版时
$18{:}00$ 值 $52.28$~万元、$6{:}00$ 值 $35.64$~万元、$12{:}00$ 值 $28.63$~万元，
且\textbf{边际价值强烈依赖已有的信息集}——只要已有 $18{:}00$，再订 $12{:}00$
\textbf{反而更贵}，而 $6{:}00$ 与 $18{:}00$ 在已有其余两版时仍分别值 $14.07$
万元与 $11.71$~万元。可见\textbf{$12{:}00$ 与 $18{:}00$ 携带的信息高度重叠，
而 $6{:}00$ 携带的是另一类信息}，故只多订一版应选 $18{:}00$。这与问题三在固定
电价下的结论方向一致，说明该重复性来自\textbf{两个时刻自身的信息结构}，而非
价格是否波动。十二条覆盖边的完整增量、"增量为负应读作\textbf{信息已无增量}
而非\textbf{多一版有害}"这一辨析及其费用构成规律，见附录第
\ref{app:p4-combo-detail} 节。

\subsubsection{模型的边界}
\label{sec:p4-limit}

本文的模型有若干明确的边界，以免读者把结论推广到未被验证的范围。其中最关键
的一条是\textbf{价格预测的能力边界}：附件 4 与附件 1 的年平均完全相同，二者
之差是一个与附件 1 曲线无关的残差，该残差中只有 $41.7\%$ 的方差可由星期结构
解释，其余部分本文按不可预测处理。这意味着\textbf{若现实中价格包含真正的随机
成分（如机组非计划停运、极端天气导致的负荷尖峰），本文的预测器无法预知，只能
靠风险余量对冲}；$\Delta J$ 与图~\ref{fig:p4-strategy} 中"未来价格已知参考"
之间的差距，正是这一不可预测部分的价格。

另外三条边界——风险以分位数余量而非场景优化处理、计费口径存在退款与不退款的
两义、以及 $4$-$3$ 的参数搜索网格截断了最优解（$(\alpha,\rho)$ 在 $8$ 个窗口
里全部落在网格左下角，故其取值\textbf{不具可解释性}）——连同"未建模的部分"
一并见附录第 \ref{app:extra} 节。
% 本文件由 03_代码/p4_tables.py 自动生成，请勿手工修改。
% 数据源：06_支撑材料/p4_results.json

\begin{table}[htbp]
  \centering
  \caption{两个交付文件在 2025 年 2 月 1 日--12 月 31 日（334 天）的实际运行结果}
  \label{tab:p4-branch}
  \footnotesize
  \begin{tabular}{lrr}
    \toprule
    指标 & result4-2（计划不得修改） & result4-3（6/12/18 可调整） \\
    \midrule
    计划量（$0{:}00$ 定下的量）/kWh & 21,800,593 & 20,514,619 \\
    最终生效常规量/kWh & 21,800,593 & 21,111,818 \\
    上调量/kWh & --- & 597,199 \\
    紧急购电量/kWh & 105,603 & 52,149 \\
    \midrule
    弃电量/kWh & 2,753,256 & 1,941,508 \\
    计划费/元 & 14,119,897 & 13,187,623 \\
    调整费/元 & --- & 714,424 \\
    紧急费/元 & 643,263 & 291,979 \\
    合计费用/元 & \textbf{14,763,160} & \textbf{14,194,026} \\
    期初储电量/kWh & 7,283 & 6,534 \\
    期末储电量/kWh & 6,102 & 5,941 \\
    \bottomrule
  \end{tabular}
  \par\vspace{2pt}
  \begin{minipage}{.92\textwidth}\footnotesize
    注：两个分支都在\textbf{同一条附件 4 实际价格路径}上、按各自的规则独立回放并独立标定，费用一栏不可相加。两者都按实际电价结算，包括初始计划费；主计费口径下不允许下调，故 result4-3 的下调量为零。期末储电量用于说明跨日状态差异——只比费用而不比期末储量是不完整的。
  \end{minipage}
\end{table}

